\chapter{Wybór, walidacja i regularizacja modeli statystycznych}
\label{ch:model-selection}
\paragraph{Powiązania i odsyłacze.} Modele bazowe pochodzą z rozdziałów~\ref{ch:linear-regression} i~\ref{ch:glm}; o ryzyku predykcji i funkcjach straty zobacz rozdział~\ref{ch:decision}.


\section{Dopasowanie nie jest jeszcze uogólnianiem}
W rozdziałach~\ref{ch:linear-regression} i~\ref{ch:glm} nauczyliśmy się
estymować współczynniki modelu oraz wnioskować o nich \emph{pod założeniem
poprawności modelu}. Praktyczny problem jest trudniejszy: czy rozważany
model dobrze przewiduje nowe dane, czy nie dopasował jedynie przypadkowego
szumu? Jeśli dodamy do regresji liniowej kolejne kolumny, resztowa suma
kwadratów na danych uczących nie wzrośnie, ale jakość predykcji może się
pogorszyć. Odróżnimy trzy cele: \emph{opis zależności}, \emph{predykcję}
i \emph{wnioskowanie przyczynowe}. Dla każdego z nich dobór modelu może być
inny. W szczególności niska wartość błędu predykcji sama nie uzasadnia
interpretacji współczynnika jako efektu przyczynowego.

\begin{definicja}[Ryzyko predykcyjne i ryzyko empiryczne]
\label{def:predictive-risk}
Niech $(\mathsf X,\mathcal X)$ oraz $(\mathsf Y,\mathcal Y)$
będą przestrzeniami mierzalnymi, $(X,Y)$ parą losową
o prawie $P$ na ich iloczynie, a
$\ell:\mathsf Y\times\mathsf Y\to[0,\infty)$ mierzalną funkcją
straty. Dla mierzalnej funkcji predykcyjnej
$f:\mathsf X\to\mathsf Y$ definiujemy
$R_P(f)=\E_P\ell(Y,f(X))$, zakładając skończoność wartości oczekiwanej.
Dla obserwacji $(X_i,Y_i)_{i=1}^n$ ryzyko empiryczne wynosi
$\widehat R_n(f)=n^{-1}\sum_{i=1}^n\ell(Y_i,f(X_i))$.
Jeżeli $\widehat f_n$ jest wyznaczona z tej samej próby, to
$\widehat R_n(\widehat f_n)$ nazywamy \emph{błędem treningowym}.
Dla niezależnej od próby uczącej pary $(X_*,Y_*)\sim P$ błąd
nowej obserwacji ma warunkową wartość oczekiwaną
$R_P(\widehat f_n)=\E[\ell(Y_*,\widehat f_n(X_*))\mid\mathcal D_n]$,
gdzie $\mathcal D_n$ oznacza dane uczące.
\end{definicja}

\begin{uwaga}[Co jest losowe?]
$R_P(f)$ jest liczbą dla ustalonego $f$; $R_P(\widehat f_n)$ jest
zmienną losową, ponieważ procedura treningowa używa losowej próby.
Błąd na zbiorze testowym estymuje ryzyko \emph{już wytrenowanego}
modelu, pod warunkiem że testowe dane reprezentują docelowy rozkład
i nie były wykorzystywane przy wyborze tego modelu.
\end{uwaga}

\begin{twierdzenie}[Optymizm błędu treningowego w normalnym modelu liniowym]
\label{thm:training-optimism}
Niech $n>p\ge1$, $X\in\R^{n\times p}$ ma rząd $p$ oraz
$Y=X\beta+\varepsilon$, gdzie
$\varepsilon\sim\mathcal N_n(0,\sigma^2I_n)$ i $\sigma>0$.
Niech $H=X(X^\top X)^{-1}X^\top$ oraz $\widehat Y=HY$.
Niech niezależny wektor $Y^{\rm new}=X\beta+\varepsilon^{\rm new}$
ma ten sam rozkład co $Y$. Wtedy
\[
 \E\|Y-\widehat Y\|_2^2=(n-p)\sigma^2,
 \qquad
 \E\|Y^{\rm new}-\widehat Y\|_2^2=(n+p)\sigma^2.
\]
Zatem optymizm błędu treningowego wynosi dokładnie
$2p\sigma^2/n$ po podzieleniu sum kwadratów przez $n$.
\end{twierdzenie}
\begin{proof}
Ponieważ $(I_n-H)X\beta=0$, mamy
$Y-\widehat Y=(I_n-H)\varepsilon$. Stąd
$\E\|Y-\widehat Y\|^2=\sigma^2\operatorname{tr}(I_n-H)
=(n-p)\sigma^2$.
Z kolei $Y^{\rm new}-\widehat Y=\varepsilon^{\rm new}-H\varepsilon$.
Składniki są niezależne, wycentrowane, więc średni kwadrat normy to
$\sigma^2\operatorname{tr}(I_n)+\sigma^2\operatorname{tr}(H^\top H)
=(n+p)\sigma^2$, gdyż $H$ jest rzutem rzędu $p$.
\end{proof}

\begin{przyklad}[Więcej parametrów nie musi pomagać]
Dwa modele liniowe o wymiarach $p=2$ i $p=12$ dopasowano do tej samej
próby $n=50$ przy poprawności obu modeli i tej samej wariancji błędu.
Różnica oczekiwanego optymizmu, liczona na obserwację, wynosi
$2(12-2)\sigma^2/50=0{,}4\sigma^2$. To nie mówi, który model
jest bezwzględnie lepszy: model prostszy może pomijać ważną
zależność, a model bogatszy zwiększa wariancję predykcji.
\end{przyklad}

\section{AIC i BIC: dwie kary za złożoność}
\begin{definicja}[Kryteria informacyjne]
Niech $\{P_\theta:\theta\in\Theta\subset\R^k\}$ będzie
regularnym modelem parametrycznym o log-wiarygodności $\ell_n(\theta)$
dla \emph{tych samych obserwacji} $y_1,\ldots,y_n$.
Jeżeli istnieje maksimum $\widehat\theta$, definiujemy
\[
 \operatorname{AIC}=-2\ell_n(\widehat\theta)+2k,
 \qquad
 \operatorname{BIC}=-2\ell_n(\widehat\theta)+k\log n.
\]
AIC to \emph{kryterium informacyjne Akaikego} (Akaike information
criterion), a BIC to \emph{bayesowskie kryterium informacyjne}
(Bayesian information criterion). W obu przypadkach mniejsza
wartość oznacza preferencję kryterium wśród porównywanych modeli.
Liczba $k$ obejmuje wszystkie
estymowane parametry, także parametr wariancji, jeżeli jest nieznany.
\end{definicja}

\begin{twierdzenie}[AIC i BIC w regresji normalnej]
\label{thm:aic-gaussian}
Niech $Y\sim\mathcal N_n(X\beta,\sigma^2I_n)$, gdzie
$X\in\R^{n\times p}$ ma pełny rząd kolumnowy, $n>p$,
a $\beta\in\R^p$ i $\sigma^2>0$ są nieznane.
Jeśli $\mathrm{RSS}=\min_b\|Y-Xb\|^2>0$, to MLE wariancji
wynosi $\widehat\sigma^2=\mathrm{RSS}/n$ i
\[
 -2\ell_n(\widehat\beta,\widehat\sigma^2)
 =n\bigl(\log(2\pi)+1+\log(\mathrm{RSS}/n)\bigr).
\]
W szczególności, przy porównywaniu modeli dla tego samego $Y$,
$\operatorname{AIC}=n\log(\mathrm{RSS}/n)+2(p+1)+C_n$ oraz
$\operatorname{BIC}=n\log(\mathrm{RSS}/n)+(p+1)\log n+C_n$,
gdzie $C_n=n(\log(2\pi)+1)$ nie zależy od modelu.
\end{twierdzenie}
\begin{proof}
Log-wiarygodność wynosi
$\ell_n(b,v)=-\tfrac n2\log(2\pi v)-\|Y-Xb\|^2/(2v)$.
Dla każdego $v>0$ maksymalizacja po $b$ prowadzi do MNK.
Dla $v>0$ funkcja $-\tfrac n2\log v-\mathrm{RSS}/(2v)$
ma pochodną $(\mathrm{RSS}-nv)/(2v^2)$ i osiąga
jedyne maksimum w $v=\mathrm{RSS}/n$, pod warunkiem
$\mathrm{RSS}>0$. Podstawienie i policzenie parametrów daje tezę.
\end{proof}

\begin{uwaga}[Dlaczego kryteria różnią się celem?]
AIC ma uzasadnienie asymptotyczne w szacowaniu pozapróbkowego
błędu logarytmicznego (straty związanej z dywergencją
Kullbacka--Leiblera). BIC powstaje, przy dodatkowych
założeniach, z przybliżenia logarytmu \emph{zintegrowanej
wiarygodności} i silniej karze liczbę parametrów, gdy $n$
jest duże. Nie są to równoważne cele. Nie należy
interpretować $\Delta\mathrm{AIC}$ jako p-value ani niższego BIC
jako dowodu prawdziwości modelu.
\end{uwaga}

\begin{stwierdzenie}[Źródło kary BIC: przybliżenie Laplace'a]
\label{prop:bic-laplace}
Niech $\Theta\subset\R^k$ będzie otwarta, a $\pi$
gęstością a priori dodatnią i ciągłą w otoczeniu
$\theta_0$. Załóżmy, że log-wiarygodność $\ell_n$
ma wewnętrzny maksymalizator $\widehat\theta_n$ zbieżny do
$\theta_0$ według prawdopodobieństwa, że lokalnie
$-n^{-1}\nabla^2\ell_n(\theta)\to J$ jednostajnie według
prawdopodobieństwa, gdzie $J$ jest dodatnio określona,
i że pozostała część całki jest zaniedbywalna względem
wkładu otoczenia maksimum. Zakładamy również
odpowiednią dominację umożliwiającą asymptotykę Laplace'a.
Wtedy dla $m_n=\int_\Theta e^{\ell_n(\theta)}\pi(\theta)
\,\dd\lambda_k(\theta)$ zachodzi
$-2\log m_n=-2\ell_n(\widehat\theta_n)+k\log n
+O_{\mathbb P}(1)$.
\end{stwierdzenie}
\begin{proof}
W otoczeniu maksimum rozwinięcie Taylora daje
$\ell_n(\theta)=\ell_n(\widehat\theta_n)-\tfrac n2
(\theta-\widehat\theta_n)^\top J
(\theta-\widehat\theta_n)+o(n\|\theta-\widehat\theta_n\|^2)$
na skali $\|\theta-\widehat\theta_n\|=O(n^{-1/2})$,
w sensie prawdopodobieństwa.
Po podstawieniu $u=\sqrt n(\theta-\widehat\theta_n)$
całka lokalna jest asymptotycznie równa
$e^{\ell_n(\widehat\theta_n)}\pi(\theta_0)n^{-k/2}
(2\pi)^{k/2}(\det J)^{-1/2}$.
Założenie o zaniedbywalności reszty pozwala zastąpić
całą całkę wkładem lokalnym; logarytmowanie kończy dowód.
\end{proof}

\begin{uwaga}[Uczciwe granice twierdzenia]
Do powyższego przybliżenia nie wystarcza samo istnienie
hesjanu w maksimum: potrzeba kontroli całki poza nim,
regularności skali $n^{-1/2}$ i właściwego prioru.
Przy parametrach na brzegu, w modelach osobliwych
lub przy silnej nieidentyfikowalności kara $k\log n$
może nie być właściwa. W szczególności pełne
uzasadnienie AIC w modelu nieregularnym wymaga
osobnej analizy; prosty wzór z karą $2k$ nie jest
uniwersalną korektą optymizmu.
\end{uwaga}

\section{Rozdzielenie danych, walidacja krzyżowa i wyciek informacji}
\begin{definicja}[Walidacja $K$-krotna]
Niech $(Z_i)_{i=1}^n$, $Z_i=(X_i,Y_i)$, będzie próbą
niezależną o wspólnym rozkładzie. Dzielimy indeksy
$\{1,\ldots,n\}$ na niepuste, rozłączne zbiory
$V_1,\ldots,V_K$, $2\le K\le n$, o sumie równej
całemu zbiorowi. Dla algorytmu uczącego $\mathcal A$
niech $\widehat f^{(-j)}$ oznacza model wytrenowany
wyłącznie na danych o indeksach spoza $V_j$.
Wtedy \emph{błąd walidacji krzyżowej} jest równy
\[
 \widehat R_{\rm CV}=
 \frac1n\sum_{j=1}^K\sum_{i\in V_j}
 \ell(Y_i,\widehat f^{(-j)}(X_i)).
\]
\end{definicja}

\begin{stwierdzenie}[Bezstronność błędu testowego przy niezależnym teście]
\label{prop:test-unbiased}
Niech $(\mathsf X,\mathcal X)$ i $(\mathsf Y,\mathcal Y)$ będą przestrzeniami mierzalnymi, $P$ prawem na $(\mathsf X\times\mathsf Y,\mathcal X\otimes\mathcal Y)$, a $m\ge1$. Niech $\mathcal D_{\rm train}$ oznacza dane uczące, $\widehat f=\mathcal A(\mathcal D_{\rm train})$ będzie mierzalną funkcją predykcyjną $\mathsf X\to\mathsf Y$, a $(X_j^*,Y_j^*)_{j=1}^m$ niezależną próbą testową o wspólnym prawie $P$, niezależną od $\mathcal D_{\rm train}$. Niech strata $\ell:\mathsf Y\times\mathsf Y\to\R$ będzie mierzalna i $\E|\ell(Y_j^*,\widehat f(X_j^*))|<\infty$. Definiujemy losowe ryzyko $R_P(\widehat f)=\int_{\mathsf X\times\mathsf Y}\ell(y,\widehat f(x))\,\dd P(x,y)$. Wówczas estymator
$\widehat R_{\rm test}=m^{-1}\sum_{j=1}^m
\ell(Y_j^*,\widehat f(X_j^*))$
spełnia $\E(\widehat R_{\rm test}\mid\mathcal D_{\rm train})
=R_P(\widehat f)$.
\end{stwierdzenie}
\begin{proof}
Po warunkowaniu względem danych uczących $\widehat f$
jest ustalona, zaś każda para $(X_j^*,Y_j^*)$ ma rozkład $P$.
Liniowość warunkowej wartości oczekiwanej daje
$\E[\ell(Y_j^*,\widehat f(X_j^*))\mid\mathcal D_{\rm train}]
=\int_{\mathsf X\times\mathsf Y}\ell(y,\widehat f(x))\,\dd P(x,y)$.
Uśrednienie po $j$ kończy dowód.
\end{proof}

\begin{przyklad}[Wyciek informacji przed walidacją]
Jeśli standaryzujemy kolumny całej tabeli, a dopiero potem
dzielimy dane na trening i test, średnie i odchylenia użyte
do transformacji korzystają z obserwacji testowych. Podobny
problem pojawia się przy imputacji braków danych, wyborze cech
na podstawie korelacji z wynikiem albo selekcji hiperparametrów
na zbiorze testowym. Każda taka operacja musi zostać wykonana
\emph{wewnątrz} odpowiedniego fragmentu uczącego.
Dla próby niezależnej służy do tego \texttt{Pipeline}; przy
selekcji hiperparametrów i równoczesnym szacowaniu jakości
stosujemy zagnieżdżoną walidację krzyżową albo oddzielny test.
\end{przyklad}

\begin{center}
\begin{tikzpicture}[>=Latex, every node/.style={font=\small},
  s/.style={draw,rounded corners,align=center,minimum height=9mm,inner sep=3pt},
  node distance=11mm and 8mm]
\node[s] (d) {Dane dostępne do analizy};
\node[s,right=of d] (tr) {Trening i walidacja\\strojenie modelu};
\node[s,right=of tr] (te) {Odłożony test\\jednorazowa ocena};
\draw[->,thick] (d)--(tr);
\draw[->,thick] (tr)--(te);
\node[below=6mm of tr,align=center,text=black!70] {Wszystkie transformacje i dobór parametrów\\wyłącznie w obrębie danych treningowych};
\end{tikzpicture}
\end{center}

\begin{uwaga}[Czego walidacja $K$-krotna nie obiecuje?]
$\widehat R_{\rm CV}$ nie jest na ogół dokładnie
nieobciążonym oszacowaniem ryzyka modelu trenowanego
na całych $n$ obserwacjach: w kolejnych podziałach
uczymy na mniejszych próbach. Przy danych zależnych
losowe mieszanie obserwacji może dodatkowo zaniżać błąd.
Dla szeregów czasowych stosujemy podziały zachowujące
chronologię, opisane w rozdziale~\ref{ch:time-series}.
\end{uwaga}

\section{Ridge i lasso: kiedy wprowadzenie obciążenia pomaga predykcji?}
\begin{definicja}[Regresja ridge i lasso]
Dla ustalonej macierzy $X\in\R^{n\times p}$, wektora
$y\in\R^n$ i $\lambda>0$ definiujemy:
\begin{itemize}
\item \emph{ridge}: dowolny minimizator
$Q_{\rm R}(b)=\frac12\|y-Xb\|_2^2+\frac\lambda2\|b\|_2^2$;
\item \emph{lasso}: dowolny minimizator
$Q_{\rm L}(b)=\frac12\|y-Xb\|_2^2+\lambda\sum_{j=1}^p|b_j|$.
\end{itemize}
W zastosowaniach z wyrazem wolnym zwykle \emph{nie karzemy}
tego wyrazu i standaryzujemy cechy liczbowe, ponieważ
wielkość kary zależy od jednostek współczynników.
\end{definicja}

\begin{twierdzenie}[Rozwiązanie ridge i jego obciążenie]
\label{thm:ridge}
Dla dowolnego $X\in\R^{n\times p}$, $y\in\R^n$
i $\lambda>0$ problem ridge ma dokładnie jedno rozwiązanie:
$\widehat\beta_\lambda=(X^\top X+\lambda I_p)^{-1}X^\top y$.
Jeżeli $Y=X\beta+\varepsilon$ z $\E\varepsilon=0$
i $\Cov(\varepsilon)=\sigma^2I_n$, to
\[
 \E\widehat\beta_\lambda=(X^\top X+\lambda I_p)^{-1}
 X^\top X\beta,
 \qquad
 \Cov(\widehat\beta_\lambda)=\sigma^2A_\lambda X^\top X A_\lambda,
 \quad A_\lambda=(X^\top X+\lambda I_p)^{-1}.
\]
\end{twierdzenie}
\begin{proof}
Hesjan funkcji $Q_{\rm R}$ wynosi
$X^\top X+\lambda I_p$, jest dodatnio określony, więc
kryterium jest ściśle wypukłe i koercywne.
Jego pochodna $X^\top(Xb-y)+\lambda b$ zeruje się
wyłącznie w podanym punkcie. Ostatnie równości
wynikają z liniowości estymatora i tożsamości
$\widehat\beta_\lambda=A_\lambda X^\top(X\beta+\varepsilon)$.
\end{proof}

\begin{twierdzenie}[Lasso z ortogonalnymi kolumnami]
\label{thm:lasso-soft}
Jeżeli $X\in\R^{n\times p}$ spełnia $X^\top X=I_p$,
a $z=X^\top y$, to jedyny minimizator $Q_{\rm L}$ ma
współrzędne
$\widehat\beta_{\lambda,j}=\operatorname{sgn}(z_j)
\max(|z_j|-\lambda,0)$, przy $\operatorname{sgn}(0)=0$.
\end{twierdzenie}
\begin{proof}
Po odrzuceniu składnika niezależnego od $b$ kryterium
jest sumą funkcji $q_j(t)=\frac12(t-z_j)^2+\lambda|t|$.
Dla $t>0$ pochodna to $t-z_j+\lambda$, dla $t<0$
wynosi $t-z_j-\lambda$. W punkcie $0$ warunek minimum
dla funkcji wypukłej to $0\in-z_j+\lambda[-1,1]$.
Z tych trzech przypadków otrzymujemy formułę.
Ścisła wypukłość składnika kwadratowego zapewnia
jedyność rozwiązania.
\end{proof}

\begin{uwaga}[Wybór ridge czy lasso]
Ridge zmniejsza współczynniki w sposób ciągły, ale na ogół
nie zeruje ich dokładnie; jest szczególnie pomocny przy
współliniowości i wielu słabych sygnałach.
Lasso może dawać współczynniki równe zeru, jednak przy
silnej korelacji cech wybór konkretnych zmiennych bywa
niestabilny. Żadna z metod nie naprawia błędnej definicji
jednostki obserwacji lub zależności czasowej. Wybór
$\lambda$ powinien korzystać wyłącznie z danych uczących.
\end{uwaga}

\section{Laboratorium: dobór stopnia wielomianu bez przecieku danych}
W kolejnym programie używamy wielomianowych regresorów
o stopniach od $1$ do $8$. Każdy kandydat ma oddzielny
potok transformacji: rozwinięcie cech i dopasowanie
regresji zachodzą tylko na danej części treningowej.
Na końcu porównujemy wybór walidacyjny z wynikiem na
odłożonym teście. W pokazanej, odtwarzalnej symulacji
najlepiej wypada stopień $2$, z walidacyjnym MSE około
$0{,}319$ i testowym MSE około $0{,}189$. Nie oczekujemy,
że wartości te będą identyczne: oba oszacowania opierają
się na różnych fragmentach skończonej próby.
Poniższy przykład jest symulacją
odtwarzalną po ustaleniu ziarna generatora; rzeczywiste
dane powinny być dzielone zgodnie z mechanizmem ich
powstawania.

\begin{lstlisting}[language=Python,caption={Kompletny program: wybór stopnia modelu, walidacja i niezależny test (Python, scikit-learn).}]
import numpy as np
from sklearn.model_selection import KFold, cross_val_score
from sklearn.pipeline import make_pipeline
from sklearn.preprocessing import PolynomialFeatures, StandardScaler
from sklearn.linear_model import Ridge
from sklearn.metrics import mean_squared_error

rng = np.random.default_rng(2026)
x = rng.uniform(-2.0, 2.0, size=150)
y = 1.5 + 0.8*x - 0.7*x**2 + rng.normal(0, 0.5, 150)
indices = rng.permutation(len(x))
train, test = indices[:110], indices[110:]
Xtr, Xte = x[train, None], x[test, None]
ytr, yte = y[train], y[test]
cv = KFold(n_splits=5, shuffle=True, random_state=11)

results = []
for degree in range(1, 9):
    pipe = make_pipeline(PolynomialFeatures(degree, include_bias=False),
                         StandardScaler(), Ridge(alpha=1.0))
    # Negative score because scikit-learn maximizes its scoring metric.
    mse_cv = -cross_val_score(pipe, Xtr, ytr, cv=cv,
                             scoring="neg_mean_squared_error").mean()
    results.append((degree, mse_cv))

best_degree, best_cv = min(results, key=lambda item: item[1])
best = make_pipeline(PolynomialFeatures(best_degree, include_bias=False),
                     StandardScaler(), Ridge(alpha=1.0))
best.fit(Xtr, ytr)
test_mse = mean_squared_error(yte, best.predict(Xte))
print("CV results:", results)
print("degree chosen by CV:", best_degree)
print("CV MSE:", best_cv)
print("holdout MSE:", test_mse)
\end{lstlisting}

\paragraph{Interpretacja.}
Program nie wybiera stopnia na podstawie wyniku testowego.
Nie porównujemy wielu konfiguracji na teście, a następnie
nie nazywamy najmniejszego błędu \emph{niezależną oceną}.
Przy rzeczywistych danych czasowych zwykły \texttt{KFold}
z mieszaniem byłby niewłaściwy; użyjemy innego podziału
w rozdziale~\ref{ch:time-series}.

\section{Porównania metod i zadania}
\begin{center}\small
\begin{tabular}{>{\raggedright\arraybackslash}p{.24\textwidth}>{\raggedright\arraybackslash}p{.32\textwidth}>{\raggedright\arraybackslash}p{.31\textwidth}}
\toprule
\textbf{Cel} & \textbf{Metoda} & \textbf{Warunek i ryzyko}\\
\midrule
Dopasowanie porównywalnych modeli o pełnej wiarygodności & AIC & Te same dane i odpowiedź; regularna asymptotyka; ryzyko przetrenowania.\\
Preferowanie prostszego modelu parametrycznego & BIC & Regularność i interpretacja w kategoriach aproksymacji ewidencji modelu.\\
Wybór jakości predykcji dla danych iid & Walidacja krzyżowa & Transformacje w ramach foldów; dobór liczby foldów i metryki.\\
Niezależna ocena końcowego modelu & Odłożony zbiór testowy & Test nie może uczestniczyć w strojeniu.\\
Wiele skorelowanych zmiennych & Ridge & Obciążenie za mniejszą wariancję; skala cech.\\
Szukamy rzadkiego zbioru predyktorów & Lasso & Niestabilna selekcja przy współliniowości.\\
\bottomrule
\end{tabular}
\end{center}

\begin{zadanie}[Dlaczego nie można porównać AIC wprost?]
Analityk dopasował regresję logistyczną do $100$ produktów,
z których $5$ ma braki danych, oraz regresję Poissona do
$95$ obserwacji dotyczących innego rodzaju odpowiedzi.
Pierwszy model ma AIC równe $89$, drugi $76$. Czy drugi
model jest lepszy?
\par\emph{Rozwiązanie.} Nie. Wiarygodności dotyczą innych
zmiennych i potencjalnie różnych zbiorów danych;
różnice AIC nie odpowiadają wspólnej mierze ryzyka.
Najpierw trzeba ustalić ten sam cel predykcji, rozkład
odpowiedzi, obserwacje i stratę. Dla dwóch konkurencyjnych
modeli tego samego wyniku można zastosować wspólną
walidację na identycznych podziałach danych.
\end{zadanie}

\begin{zadanie}[Ridge a nieobciążoność]
Niech $X^\top X=I_p$, a $Y=X\beta+\varepsilon$,
$\E\varepsilon=0$, $\Cov(\varepsilon)=\sigma^2I_n$.
Wyznacz obciążenie ridge oraz MSE wektorowy
$\E\|\widehat\beta_\lambda-\beta\|_2^2$.
\par\emph{Rozwiązanie.} Z twierdzenia~\ref{thm:ridge}
$\E\widehat\beta_\lambda=\beta/(1+\lambda)$,
więc obciążenie wynosi $-\lambda\beta/(1+\lambda)$.
Kowariancja jest równa
$\sigma^2(1+\lambda)^{-2}I_p$, skąd
$\E\|\widehat\beta_\lambda-\beta\|_2^2
=(\lambda^2\|\beta\|_2^2+p\sigma^2)/(1+\lambda)^2$.
Zależnie od nieznanego $\beta$ pewne dodatnie
$\lambda$ mogą poprawiać MSE względem MNK, ale nie
wynika z tego uniwersalna optymalność jednej kary.
\end{zadanie}

% ================================================================
% ROZDZIAL 16 -- DANE ZALEZNE W CZASIE
% ================================================================
